\chapter{System Design e buone pratiche di progettazione}
% Il codice in linea (\code) non va a capo: un po' di elasticita' in piu' evita righe che sbordano
\emergencystretch=2.5em

Con i requisiti e i modelli di analisi abbiamo descritto \textbf{che cosa} il sistema deve fare. Da qui in avanti ci occupiamo
di \textbf{come} farlo. Questo capitolo raccoglie le lezioni dedicate al \textit{System Design} e segue due fili:
\begin{itemize}
  \item le \textbf{attività} della progettazione di sistema: fissare gli obiettivi di progetto e scomporre il sistema in
        sottosistemi;
  \item le \textbf{regole di buona progettazione}, valide per qualunque software orientato agli oggetti: coesione,
        accoppiamento, Legge di Demetra e progettazione guidata dalle responsabilità.
\end{itemize}
La fase era già stata presentata a grandi linee nel Capitolo 3 (\textit{Fase 2: System Design}); qui la riprendiamo in
dettaglio.

% ============================================================
\section{Dove siamo nel processo}

\waterfallbar{2}

Nel modello a cascata il System Design segue la Requirements Engineering. In ingresso riceve i requisiti, i casi d'uso e il
modello di dominio; in uscita produce l'\textbf{architettura del sistema}. I requisiti vengono assegnati ai sottosistemi
software, i sottosistemi alle risorse hardware, e si scelgono i \textbf{pattern architetturali} da adottare.

\dfn{System Design}{
  Il \term{System Design} è l'insieme delle attività con cui l'ingegnere del software trasforma il \textbf{modello di analisi}
  nel \textbf{modello di progetto del sistema}.

  L'obiettivo finale del System Design e del successivo Object Design è produrre un nuovo documento, scritto in linguaggio
  \textbf{tecnico e formale}, da consegnare ai programmatori perché possano implementare il software descritto nel documento
  dei requisiti.
}

La progettazione si svolge su due livelli, che corrispondono alle fasi 2 e 3 della cascata.

\Needspace{14\baselineskip}
\begin{center}
\begin{tikzpicture}
  \node[card=ph2, text width=6.6cm, align=left, anchor=north] (sd) at (0,0)
    {\textbf{\color{ph2!80!black}System Design}\\[3pt]
     \begin{enumerate}[leftmargin=1.2em, itemsep=1pt, topsep=1pt]
       \item definire gli \textbf{obiettivi di progetto};
       \item definire l'\textbf{architettura}, scomponendo il sistema in sottosistemi, che team diversi possono implementare
             \textbf{in parallelo};
       \item scegliere le \textbf{strategie} di costruzione: hardware e software, gestione dei dati persistenti, flusso di
             controllo globale, politiche di controllo degli accessi\ldots
     \end{enumerate}};
  \node[card=ph3, text width=6.6cm, align=left, anchor=north] (od) at (7.6,0)
    {\textbf{\color{ph3!80!black}Object Design}\\[3pt]
     \begin{itemize}[leftmargin=1.2em, itemsep=1pt, topsep=1pt]
       \item data l'architettura, specificare i \textbf{dettagli implementativi} dei singoli sottosistemi;
       \item scegliere le strategie di implementazione migliori, ad esempio i \textbf{design pattern}.
     \end{itemize}};
  \draw[flow] (sd.east |- 0,-1.2) -- (od.west |- 0,-1.2);
\end{tikzpicture}
\end{center}

\Needspace{16\baselineskip}
\subsection{Cambia l'interlocutore}

Con la progettazione il punto di vista cambia radicalmente. Finora abbiamo lavorato \textbf{per il cliente} e con il
cliente, usando un linguaggio da non addetti ai lavori. Nel System Design i nostri interlocutori sono i \textbf{team di
programmatori} che dovranno implementare il sistema: chi legge i nostri documenti è un esperto di informatica, e il
linguaggio diventa tecnico.

\begin{center}
\begin{tikzpicture}
  \node[card=ph1, text width=5.4cm] (a) at (0,0)
    {\textbf{Analisi}\\[2pt] parliamo con il \textbf{cliente}\\ linguaggio per profani, poco formale\\
     \textit{casi d'uso, scenari, mock-up}};
  \node[card=ph2, text width=5.4cm] (p) at (7.4,0)
    {\textbf{Progettazione}\\[2pt] parliamo con i \textbf{programmatori}\\ linguaggio tecnico e formale\\
     \textit{sottosistemi, interfacce, classi}};
  \draw[flow] (a) -- node[note, above]{cambio di} node[note, below]{destinatario} (p);
\end{tikzpicture}
\end{center}

% ============================================================
\section{Le attività della progettazione di sistema}

Il System Design si articola in tre macro-attività.
\begin{enumerate}
  \item \textbf{Identificare gli obiettivi di progetto.} Gli sviluppatori stabiliscono quali caratteristiche di qualità vanno
        ottimizzate e con quale priorità.
  \item \textbf{Progettare la scomposizione in sottosistemi.} A partire dai casi d'uso e dai modelli di analisi, si divide
        il sistema in parti più piccole, usando stili architetturali standard.
  \item \textbf{Raffinare la scomposizione per soddisfare gli obiettivi.} La prima scomposizione di solito non soddisfa gli
        obiettivi di progetto: la si raffina finché non li soddisfa.
\end{enumerate}

Nel dettaglio, la progettazione di sistema tocca otto aspetti, riassunti in figura. In questo capitolo trattiamo i primi
due, evidenziati, insieme alle regole generali che guidano la scomposizione.

\begin{figure}[H]
\centering
\begin{tikzpicture}[
    passo/.style={rectangle, rounded corners=3pt, draw=#1!70!black, fill=#1!8, line width=0.8pt,
                  text width=3.25cm, align=left, font=\scriptsize, inner sep=4pt, minimum height=1.55cm},
    passo/.default=swgray,
    ramo/.style={line width=0.7pt, draw=swgray!80}]
  \node[chip=swgray, font=\small\bfseries, inner sep=5pt] (c) at (5.85,0) {System Design};
  % riga superiore: passi 1-4
  \node[passo=ph2, anchor=south] (p1) at (0,1.0)
    {\textbf{1. Obiettivi di progetto}\\ definizione\\ compromessi (\textit{trade-off})};
  \node[passo=ph2, anchor=south] (p2) at (3.9,1.0)
    {\textbf{2. Scomposizione del sistema}\\ livelli e partizioni\\ coesione e accoppiamento};
  \node[passo, anchor=south] (p3) at (7.8,1.0)
    {\textbf{3. Concorrenza}\\ identificazione dei thread};
  \node[passo, anchor=south] (p4) at (11.7,1.0)
    {\textbf{4. Mapping hardware/software}\\ hardware dedicato, \textit{buy or build}\\ allocazione, connettività};
  % riga inferiore: passi 5-8
  \node[passo, anchor=north] (p5) at (0,-1.0)
    {\textbf{5. Gestione dei dati}\\ oggetti persistenti\\ file, database\\ strutture dati};
  \node[passo, anchor=north] (p6) at (3.9,-1.0)
    {\textbf{6. Risorse globali}\\ controllo degli accessi\\ sicurezza};
  \node[passo, anchor=north] (p7) at (7.8,-1.0)
    {\textbf{7. Controllo del software}\\ monolitico, a eventi\\ processi concorrenti};
  \node[passo, anchor=north] (p8) at (11.7,-1.0)
    {\textbf{8. Condizioni limite}\\ inizializzazione\\ terminazione, guasti};
  \foreach \p in {p1,p2,p3,p4} \draw[ramo] (c.north) -- (\p.south);
  \foreach \p in {p5,p6,p7,p8} \draw[ramo] (c.south) -- (\p.north);
\end{tikzpicture}
\caption{Gli otto aspetti del System Design. In questo capitolo: obiettivi di progetto e scomposizione.}
\end{figure}

% ============================================================
\section{Gli obiettivi di progetto}

Identificare le \textbf{qualità} su cui il sistema deve puntare è il primo passo della progettazione.
\begin{itemize}
  \item Molti obiettivi di progetto derivano dai \textbf{requisiti non funzionali} (Capitolo 5) o dal dominio applicativo;
        altri sono indicati direttamente dal cliente.
  \item Vanno \textbf{formalizzati esplicitamente}, perché ogni decisione di progetto importante deve essere presa in base
        allo \textbf{stesso insieme di criteri}. Se ogni sviluppatore ottimizza ciò che gli sembra più importante, il
        risultato è incoerente.
\end{itemize}

\Needspace{18\baselineskip}
\subsection{I criteri di progetto}

Gli obiettivi si scelgono da un lungo elenco di qualità desiderabili, organizzabili in cinque gruppi. Molte di queste
qualità le abbiamo già incontrate nel modello ISO/IEC 25010 del Capitolo 4: nella terza colonna indichiamo la
caratteristica corrispondente.

\begin{center}
\renewcommand{\arraystretch}{1.35}
\begin{tabular}{@{}>{\bfseries}l p{6.3cm} p{4.6cm}@{}}
  \toprule
  \normalfont\textbf{Gruppo} & \textbf{Esempi di criteri} & \textbf{In ISO/IEC 25010} \\
  \midrule
  \color{q1}Prestazioni   & tempo di risposta, throughput, memoria richiesta & efficienza prestazionale \\
  \color{q6}Affidabilità  & robustezza, disponibilità, tolleranza ai guasti, sicurezza & affidabilità, sicurezza \\
  \color{q5}Costi         & di sviluppo, di manutenzione, di gestione & (non è una qualità del prodotto) \\
  \color{q2}Manutenzione  & estendibilità, modificabilità, adattabilità, portabilità & manutenibilità, flessibilità \\
  \color{q8}Usabilità     & facilità d'uso e di apprendimento & usabilità \\
  \bottomrule
\end{tabular}
\end{center}

\subsection{I compromessi di progetto}

Nella pratica si riesce a tenere conto solo di un \textbf{piccolo sottoinsieme} di questi criteri: non è realistico, ad
esempio, sviluppare un software che sia allo stesso tempo molto sicuro e poco costoso. Gli sviluppatori devono quindi
\textbf{dare priorità} agli obiettivi, tenendo conto anche di vincoli gestionali come i tempi e il budget. I compromessi
più frequenti sono quattro.

\Needspace{20\baselineskip}
\begin{center}
\begin{tikzpicture}[
    tradeoff/.style={rectangle, rounded corners=3pt, draw=#1!70!black, fill=#1!7, line width=0.8pt,
                     text width=6.7cm, align=left, font=\scriptsize, inner sep=5pt, minimum height=2.15cm}]
  \node[tradeoff=q1, anchor=north west] (t1) at (0,0)
    {{\small\bfseries Memoria \textit{vs} velocità}\\[2pt]
     Se il software non rispetta i requisiti di tempo di risposta o throughput, si usa più memoria per renderlo veloce
     (cache, ridondanza). Se non rispetta quelli di memoria, si comprimono i dati a scapito della velocità.};
  \node[tradeoff=q2, anchor=north west] (t2) at (7.5,0)
    {{\small\bfseries Tempi di rilascio \textit{vs} funzionalità}\\[2pt]
     Se i tempi sono stretti, si può rilasciare in tempo una versione con \textbf{meno funzionalità} di quelle richieste,
     rimandando le altre.};
  \node[tradeoff=q6, anchor=north west] (t3) at (0,-2.45)
    {{\small\bfseries Tempi di rilascio \textit{vs} qualità}\\[2pt]
     Con tempi stretti, il project manager può rilasciare nei tempi un software con difetti noti e correggerli dopo,
     oppure rilasciarlo in ritardo con meno difetti.};
  \node[tradeoff=q5, anchor=north west] (t4) at (7.5,-2.45)
    {{\small\bfseries Tempi di rilascio \textit{vs} personale}\\[2pt]
     Può essere necessario \textbf{aggiungere persone} al progetto per aumentarne la produttività (ma attenzione, vedi
     il riquadro qui sotto).};
\end{tikzpicture}
\end{center}

\warning{La legge di Brooks}{
  Aggiungere persone funziona solo se lo si fa \textbf{per tempo}. Fred Brooks, in \textit{The Mythical Man-Month} (1975),
  osservò che \textit{aggiungere personale a un progetto software in ritardo lo fa ritardare ancora di più}: i nuovi arrivati
  vanno formati da chi già lavora al progetto, e il numero di canali di comunicazione cresce con il quadrato delle persone.
  Il compromesso tempi/personale va quindi valutato all'inizio, non a progetto già in ritardo.
}

\begin{esempio}{Obiettivi di progetto per la gestione delle emergenze}
Gli esempi delle slide su database e sottosistemi (\textit{Resource}, \textit{Map}, \textit{Incident Management}) vengono da
un sistema per la gestione delle emergenze: gli operatori registrano gli incidenti, consultano una mappa e assegnano le
risorse (ambulanze, vigili del fuoco). Per un sistema del genere una scelta ragionevole di priorità è:
\begin{enumerate}
  \item \textbf{affidabilità e disponibilità}: il sistema deve funzionare proprio quando c'è un'emergenza;
  \item \textbf{tempo di risposta}: un incidente registrato deve comparire subito sulla mappa degli operatori;
  \item \textbf{usabilità}: gli operatori lavorano sotto stress, l'interfaccia non deve indurli in errore.
\end{enumerate}
Il \textbf{costo} passa in secondo piano: si accetta di spendere di più, ad esempio per server ridondanti, pur di garantire
la disponibilità. È un compromesso costo/affidabilità deciso \textbf{una volta per tutte} e scritto nel documento di
progetto, in modo che ogni scelta successiva lo rispetti.
\end{esempio}

% ============================================================
\section{Criteri generali: progettare per il cambiamento}

Molti criteri di progetto dipendono strettamente dal problema da risolvere: trovare il miglior equilibrio fra i vincoli è
compito dell'analista, che si basa sulla propria esperienza e sensibilità. Esistono però anche criteri \textbf{generali},
validi per qualunque software orientato agli oggetti. Applicati correttamente, migliorano in modo significativo la
\textbf{qualità interna} del codice senza introdurre svantaggi particolari. Il resto del capitolo è dedicato a questi
criteri.

Il primo è un atteggiamento: quando progettiamo una classe dobbiamo sforzarci di pensare a quali \textbf{cambiamenti}
potranno essere richiesti in futuro.

\dfn{Progettare per anticipare il cambiamento}{
  Le scelte iniziali devono \textbf{facilitare l'evoluzione futura} del software (\textit{designing to anticipate change}).
  È uno dei requisiti fondamentali di un buon progetto, ed è anche un criterio con cui il progetto viene valutato.
}

Ricordiamo dal Capitolo 2 che la manutenzione rappresenta la parte più consistente del costo di un software: un progetto
che rende economici i cambiamenti ripaga questo investimento per tutta la vita del sistema.

% ============================================================
\section{La scomposizione in sottosistemi}

Per ridurre la complessità della soluzione si scompone il sistema in parti più piccole, i \term{sottosistemi}.
\begin{itemize}
  \item Un sottosistema corrisponde tipicamente alla porzione di lavoro che un \textbf{singolo sviluppatore o un team} può
        svolgere in modo indipendente.
  \item Se i sottosistemi sono relativamente indipendenti, i team possono lavorarci con un \textbf{costo di comunicazione}
        minimo.
  \item Se un sottosistema è ancora complesso, si applica di nuovo lo stesso principio e lo si scompone in sottosistemi più
        semplici.
\end{itemize}

\Needspace{22\baselineskip}
\subsection{Perché scomporre conviene}

Indichiamo con $P$ un problema, con $C(P)$ la sua complessità e con $E(P)$ lo sforzo necessario a risolverlo. Dati due
problemi $P_1$ e $P_2$, è intuitivo che
\[
  C(P_1) > C(P_2) \;\Rightarrow\; E(P_1) > E(P_2).
\]
Si osserva però anche che la complessità di due problemi affrontati \textbf{insieme} è maggiore della somma delle loro
complessità prese separatamente:
\[
  C(P_1 + P_2) > C(P_1) + C(P_2) \qquad\text{e quindi}\qquad E(P_1 + P_2) > E(P_1) + E(P_2).
\]

\noindent\begin{minipage}[c]{0.56\textwidth}
Il motivo sono le \textbf{interazioni}: risolvendo i due problemi insieme bisogna tenere conto anche di come ogni parte
dell'uno influenza ogni parte dell'altro. In figura, sei elementi tutti potenzialmente legati fra loro hanno
$6\cdot 5/2 = 15$ collegamenti possibili; divisi in due gruppi da tre, con un solo punto di contatto, i collegamenti sono
$3 + 3 + 1 = 7$. È l'argomento a favore del \textit{divide et impera}, a patto che i pezzi siano davvero
\textbf{indipendenti}: è quello che misureremo con coesione e accoppiamento.
\end{minipage}\hfill
\begin{minipage}[c]{0.4\textwidth}
\centering
\begin{tikzpicture}[pt/.style={circle, fill=#1, inner sep=0pt, minimum size=0.22cm}, pt/.default=swgray]
  \foreach \i in {1,...,6} \coordinate (a\i) at ({60*\i}:0.95);
  \foreach \i in {1,...,6} \foreach \j in {1,...,6} {\ifnum\i<\j \draw[swgray!60, line width=0.5pt] (a\i) -- (a\j);\fi}
  \foreach \i in {1,...,6} \node[pt=swred] at (a\i) {};
  \node[note, align=center] at (0,-1.45) {insieme: 15 legami};
  \begin{scope}[shift={(3.4,0)}]
    \coordinate (b1) at (-0.75,0.7); \coordinate (b2) at (-1.05,-0.2); \coordinate (b3) at (-0.35,-0.45);
    \coordinate (c1) at (0.75,0.7);  \coordinate (c2) at (1.05,-0.2);  \coordinate (c3) at (0.35,-0.45);
    \draw[swgreen!70, line width=0.5pt] (b1)--(b2)--(b3)--(b1) (c1)--(c2)--(c3)--(c1);
    \draw[swblue, line width=0.8pt] (b3) -- (c3);
    \foreach \p in {b1,b2,b3,c1,c2,c3} \node[pt=swgreen] at (\p) {};
    \node[note, align=center] at (0,-1.45) {scomposti: 7 legami};
  \end{scope}
\end{tikzpicture}
\end{minipage}

\subsection{Modellare i sottosistemi}

In UML un sottosistema si rappresenta con un \textbf{package} (Capitolo 8): una collezione di classi, associazioni,
operazioni e vincoli fra loro correlati. Anche i linguaggi offrono costrutti per lo stesso concetto: i \code{package} in
Java, i \code{namespace} in C++ e C\#.

\dfn{Servizi e interfaccia di un sottosistema}{
  Un sottosistema è caratterizzato dai \term{servizi} che offre agli altri sottosistemi. L'insieme delle operazioni
  disponibili agli altri sottosistemi forma l'\term{interfaccia del sottosistema}: comprende i nomi delle operazioni, i loro
  parametri, i loro tipi e i valori restituiti.
}

Il System Design si concentra sulla definizione dei servizi offerti da ogni sottosistema, in termini di
\textbf{operazioni}, \textbf{parametri} e \textbf{comportamento ad alto livello}. I dettagli interni restano all'Object
Design.

\begin{esempio}{Accesso diretto al database o tramite un sottosistema dedicato}
Torniamo al sistema per le emergenze. Nella prima versione (a sinistra) tutti i sottosistemi accedono \textbf{direttamente}
al database.
\begin{itemize}
  \item Non si riesce a specificare con chiarezza le responsabilità dei moduli: chi \textit{gestisce} i dati, chi li
        \textit{salva}?
  \item Tutti i sottosistemi dipendono dall'interfaccia del database: se cambia il DBMS, va modificato il codice di
        \textbf{tutte} le classi che vi accedono.
\end{itemize}
Nella seconda versione (a destra) si introduce il sottosistema \textbf{DataAccess}, l'unico che gestisce l'input/output con
il database. Se cambia il DBMS, si modifica \textbf{solo} DataAccess.

\begin{center}
\begin{tikzpicture}
  % versione 1
  \umlpkg{r1}{(0,2.2)}{1.9}{0.75}{\shortstack{Resource\\Management}}
  \umlpkg{m1}{(2.3,2.2)}{1.9}{0.75}{\shortstack{Map\\Management}}
  \umlpkg{i1}{(4.6,2.2)}{1.9}{0.75}{\shortstack{Incident\\Management}}
  \umlpkg{d1}{(2.3,0)}{1.9}{0.6}{Database}
  \draw[dipend, draw=swred] (r1.south) -- (d1.north west);
  \draw[dipend, draw=swred] (m1.south) -- (d1.north);
  \draw[dipend, draw=swred] (i1.south) -- (d1.north east);
  \node[note, text=swred] at (3.25,-0.45) {tre dipendenze dal DBMS};
  % versione 2
  \begin{scope}[shift={(8.2,0)}]
    \umlpkg{r2}{(0,3.2)}{1.9}{0.75}{\shortstack{Resource\\Management}}
    \umlpkg{m2}{(2.3,3.2)}{1.9}{0.75}{\shortstack{Map\\Management}}
    \umlpkg{i2}{(4.6,3.2)}{1.9}{0.75}{\shortstack{Incident\\Management}}
    \umlpkg{a2}{(2.3,1.55)}{1.9}{0.6}{DataAccess}
    \umlpkg{d2}{(2.3,0)}{1.9}{0.6}{Database}
    \draw[dipend] (r2.south) -- (a2.north west);
    \draw[dipend] (m2.south) -- (a2.north);
    \draw[dipend] (i2.south) -- (a2.north east);
    \draw[dipend, draw=swgreen, line width=0.9pt] (a2.south) -- (d2.north);
    \node[note, text=swgreen] at (3.25,-0.45) {una sola dipendenza dal DBMS};
  \end{scope}
\end{tikzpicture}
\end{center}
Nelle slide il testo dice \textit{«modifico solo il sottosistema Storage»}, ma nel diagramma lo stesso sottosistema si
chiama \textit{DataAccess}: sono la stessa cosa.
\end{esempio}

\subsection{Quanti sottosistemi?}

L'esempio ha ridotto le relazioni fra i sottosistemi, ma ha \textbf{aumentato la complessità}: c'è un sottosistema in più
da progettare, implementare e attraversare a ogni accesso ai dati. Inseguendo la riduzione delle dipendenze si rischia di
aggiungere livelli di astrazione che costano tempo di sviluppo e di esecuzione.

\Needspace{20\baselineskip}
\noindent\begin{minipage}[c]{0.44\textwidth}
Il grafico, tratto da Pressman, mostra il compromesso. Al crescere del numero di componenti:
\begin{itemize}
  \item il \textbf{costo di sviluppo} di ciascun componente diminuisce, perché ogni componente è più piccolo e semplice;
  \item il \textbf{costo di integrazione} aumenta, perché ci sono più interfacce da definire e più pezzi da far
        collaborare.
\end{itemize}
Il costo totale ha quindi un \textbf{minimo}: troppi pochi componenti sono difficili da sviluppare, troppi sono difficili
da integrare.
\end{minipage}\hfill
\begin{minipage}[c]{0.53\textwidth}
\centering
\begin{tikzpicture}[x=0.62cm, y=0.62cm]
  \foreach \n in {1,...,12} {
    \pgfmathsetmacro{\integ}{0.03*\n*\n+0.15}
    \pgfmathsetmacro{\svil}{6.2/\n+0.25}
    \fill[q6!35, draw=q6!70!black, line width=0.4pt] (\n-0.38,0) rectangle (\n+0.38,\integ);
    \fill[sworange!75, draw=q6!70!black, line width=0.4pt] (\n-0.38,\integ) rectangle (\n+0.38,\integ+\svil);
  }
  \draw[-{Stealth[length=2mm]}, line width=0.8pt] (0.3,0) -- (13.2,0) node[below left, note, align=right]
    {numero di componenti};
  \draw[-{Stealth[length=2mm]}, line width=0.8pt] (0.3,0) -- (0.3,8.2) node[above right, note] {costo del software};
  \draw[dashed, line width=0.7pt, swblue] (5,0) -- (5,4.0);
  \node[note, text=swblue, anchor=north, align=center] at (5,-0.1) {numero\\ ottimale};
  \node[note, anchor=west] at (7.4,7.4) {\textcolor{sworange!90!black}{\rule{6pt}{6pt}} costo di sviluppo};
  \node[note, anchor=west] at (7.4,6.7) {\textcolor{q6!60}{\rule{6pt}{6pt}} costo di integrazione};
\end{tikzpicture}
\end{minipage}

\subsection{L'indipendenza funzionale}

La divisione in sottosistemi è guidata da due fattori qualitativi fondamentali, la \term{coesione} e
l'\term{accoppiamento}. L'obiettivo è l'\textbf{indipendenza funzionale} dei sottosistemi, che si massimizza con
\textbf{alta coesione} e \textbf{basso accoppiamento}.

\begin{figure}[H]
\centering
\begin{tikzpicture}[
    pt/.style={circle, fill=swgray, inner sep=0pt, minimum size=0.2cm},
    mod/.style={rounded corners=5pt, draw=#1!70!black, fill=#1!6, line width=0.8pt},
    int/.style={swgreen!80!black, line width=0.6pt},
    ext/.style={swred, line width=0.7pt}]
  % --- buona scomposizione ---
  \foreach \k/\x in {A/0, B/2.6, C/5.2} {
    \draw[mod=swgreen] (\x-0.9,-0.9) rectangle (\x+0.9,0.9);
    \coordinate (\k1) at (\x-0.45,0.45); \coordinate (\k2) at (\x+0.45,0.45);
    \coordinate (\k3) at (\x-0.45,-0.45); \coordinate (\k4) at (\x+0.45,-0.45);
    \draw[int] (\k1)--(\k2)--(\k4)--(\k3)--(\k1)--(\k4) (\k2)--(\k3);
  }
  \draw[ext] (A2) to[bend left=20] (B1);
  \draw[ext] (B4) to[bend right=20] (C3);
  \foreach \k in {A,B,C} \foreach \i in {1,...,4} \node[pt] at (\k\i) {};
  \node[font=\scriptsize\bfseries, text=swgreen!70!black] at (2.6,-1.35) {alta coesione, basso accoppiamento};
  % --- cattiva scomposizione ---
  \begin{scope}[shift={(8.4,0)}]
    \foreach \k/\x in {A/0, B/2.6, C/5.2} {
      \draw[mod=swred] (\x-0.9,-0.9) rectangle (\x+0.9,0.9);
      \coordinate (\k1) at (\x-0.45,0.45); \coordinate (\k2) at (\x+0.45,0.45);
      \coordinate (\k3) at (\x-0.45,-0.45); \coordinate (\k4) at (\x+0.45,-0.45);
    }
    \draw[int] (A1)--(A3) (B2)--(B4) (C1)--(C2);
    \draw[ext] (A2) -- (B1) (A4) -- (B3) (A2) to[bend left=25] (C1) (B2) -- (C1) (B4) -- (C3)
               (A4) to[bend right=25] (C4) (B3) to[bend right=30] (C3) (A1) to[bend left=35] (B2);
    \foreach \k in {A,B,C} \foreach \i in {1,...,4} \node[pt] at (\k\i) {};
    \node[font=\scriptsize\bfseries, text=swred] at (2.6,-1.35) {bassa coesione, alto accoppiamento};
  \end{scope}
\end{tikzpicture}
\caption{Gli stessi elementi raggruppati in due modi. In verde i legami interni ai moduli, in rosso quelli fra moduli. A
         sinistra ogni modulo è quasi autosufficiente; a destra quasi ogni legame attraversa i confini.}
\end{figure}

% ============================================================
\section{Coesione}

\dfn{Coesione}{
  La \term{coesione} misura quanto sono \textbf{strettamente correlate e focalizzate} le responsabilità (o i servizi
  offerti) di un modulo. Se ogni unità è responsabile di un \textbf{unico compito}, diciamo che ha \textbf{alta coesione}.
}

L'alta coesione è una proprietà desiderabile del codice, perché permette di:
\begin{itemize}
  \item capire meglio il ruolo di una classe;
  \item riusare una classe;
  \item mantenere una classe;
  \item limitare e concentrare le modifiche al codice;
  \item usare nomi appropriati, efficaci e comunicativi.
\end{itemize}

La coesione si valuta a due livelli di \textbf{granularità}:
\begin{itemize}
  \item \textbf{coesione dei metodi}: un metodo dovrebbe essere responsabile di un solo compito ben definito;
  \item \textbf{coesione delle classi}: ogni classe dovrebbe rappresentare un solo concetto ben definito.
\end{itemize}

\subsection{Il principio di singola responsabilità}

\dfn{Single Responsibility Principle (SRP)}{
  Una classe dovrebbe avere \textbf{una sola responsabilità}, cioè \textbf{un solo motivo per cambiare}.
}

Una responsabilità è quindi un \textbf{motivo di cambiamento}: le responsabilità di una classe sono i suoi \textit{assi di
cambiamento}. Se una classe ha due responsabilità, nel progetto queste sono \textbf{accoppiate} e devono cambiare insieme.
Una classe con più di una responsabilità diventa \textbf{fragile}: si rompe ogni volta che cambia uno qualunque dei
requisiti da cui dipende.

\begin{esempio}{Un report con troppe responsabilità}
La classe \code{Report} recupera i dati, li formatta e stampa il report. L'esempio delle slide è in C\#; qui lo riportiamo
in Java.

\Needspace{16\baselineskip}
\noindent\begin{minipage}[c]{0.62\textwidth}
\begin{lstlisting}[language=Java]
public class Report {
  private List<ReportDataElement> getData() {
    System.out.println("Getting data...");
    return new ArrayList<>();
  }
  private void formatReport() {
    System.out.println("Formatting report...");
  }
  public void print() {
    getData();
    formatReport();
    System.out.println("Printing report...");
  }
}
\end{lstlisting}
\end{minipage}\hfill
\begin{minipage}[c]{0.33\textwidth}
\centering
\begin{tikzpicture}
  \umlclass[2.6cm]{rep}{(0,0)}{Report}{}{getData()\\formatReport()\\print()}
\end{tikzpicture}

\vspace{4pt}
{\scriptsize Tre motivi per cambiare:\\ la sorgente dei dati, il formato,\\ il dispositivo di stampa.\par}
\end{minipage}

Applicando l'SRP si separano le tre responsabilità in tre classi; \code{Report} si limita a chiedere la stampa.

\begin{center}
\begin{tikzpicture}
  \umlclass[2.4cm]{r}{(0,0)}{Report}{}{print()}
  \umlclass[2.4cm]{p}{(4.2,0)}{ReportPrinter}{}{print()}
  \umlclass[2.4cm]{f}{(2.4,-2.3)}{ReportFormatter}{}{formatReport()}
  \umlclass[2.4cm]{d}{(6.0,-2.3)}{DataAccess}{}{getData()}
  \draw[assoc] (r.east) -- (p.west);
  \draw[assoc] (p.south) ++(-0.5,0) |- ++(0,-0.45) -| (f.north);
  \draw[assoc] (p.south) ++(0.5,0) |- ++(0,-0.45) -| (d.north);
\end{tikzpicture}
\end{center}

\Needspace{24\baselineskip}
\noindent\begin{minipage}[t]{0.49\textwidth}
\begin{lstlisting}[language=Java]
public class DataAccess {
  public List<ReportDataElement> getData() {
    System.out.println("Getting data...");
    return new ArrayList<>();
  }
}

public class ReportFormatter {
  public void formatReport() {
    System.out.println("Formatting report...");
  }
}
\end{lstlisting}
\end{minipage}\hfill
\begin{minipage}[t]{0.49\textwidth}
\begin{lstlisting}[language=Java]
public class ReportPrinter {
  public void print() {
    DataAccess dataAccess = new DataAccess();
    ReportFormatter reportFormatter =
        new ReportFormatter();
    dataAccess.getData();
    reportFormatter.formatReport();
    System.out.println("Printing report...");
  }
}

public class Report {
  public void print() {
    ReportPrinter reportPrinter = new ReportPrinter();
    reportPrinter.print();
  }
}
\end{lstlisting}
\end{minipage}

Ora ogni motivo di cambiamento tocca una sola classe:

\begin{center}
\renewcommand{\arraystretch}{1.3}
\begin{tabular}{@{}ll@{}}
  \toprule
  \textbf{Se cambia\ldots} & \textbf{si modifica solo\ldots} \\
  \midrule
  la sorgente dei dati (un altro database, un file, un servizio web) & \code{DataAccess} \\
  il formato del report (PDF, HTML, tabella) & \code{ReportFormatter} \\
  il modo di stampare (stampante, schermo, e-mail) & \code{ReportPrinter} \\
  \bottomrule
\end{tabular}
\end{center}
\end{esempio}

\info{{Coesione sì, ma l'accoppiamento?}}{
  Nella versione divisa \code{ReportPrinter} crea con \code{new} le classi \textbf{concrete} \code{DataAccess} e
  \code{ReportFormatter}: per stampare un report letto da un file invece che da un database bisognerebbe comunque
  modificare \code{ReportPrinter}. La coesione è migliorata, ma l'accoppiamento resta alto. Vedremo nella sezione
  sull'accoppiamento come ridurlo con le \textbf{interfacce}.

  Si noti anche che nel codice delle slide i dati letti da \code{getData()} non vengono passati a \code{formatReport()}:
  è uno schema che si limita a stampare messaggi. In una versione reale i dati scorrerebbero da \code{DataAccess} a
  \code{ReportFormatter}, ad esempio con \code{formatReport(dati)}.
}

\subsection{I vantaggi della coesione}

\Needspace{8\baselineskip}
\begin{itemize}
  \item \textbf{Riuso}: se tutte le classi rispettano l'SRP, è più probabile riusarne qualcuna.
  \item \textbf{Chiarezza}: il codice è più pulito, perché le classi non fanno cose inattese.
  \item \textbf{Nomi}: se una classe ha una sola responsabilità, trovarle un buon nome è facile.
  \item \textbf{Leggibilità}: chiarezza, nomi migliori e file più corti migliorano molto la leggibilità.
\end{itemize}

\Needspace{14\baselineskip}
\noindent\begin{minipage}[t]{0.49\textwidth}
\gbox{Indicatori di alta coesione}{swgreen}{\raggedright
\begin{itemize}[leftmargin=1.1em]
  \item la classe ha responsabilità moderate, limitate a una sola area funzionale;
  \item collabora con altre classi per completare i compiti;
  \item ha un numero limitato di metodi, cioè di funzionalità fortemente correlate;
  \item tutti i metodi sembrano appartenere allo stesso insieme, con uno scopo complessivo simile;
  \item la classe è facile da capire.
\end{itemize}}
\end{minipage}\hfill
\begin{minipage}[t]{0.49\textwidth}
\gbox{Campanelli d'allarme}{swred}{\raggedright
\begin{itemize}[leftmargin=1.1em]
  \item il nome contiene «e» o termini generici (\code{GestoreUtentiEStampe}, \code{Manager}, \code{Utils});
  \item gruppi di metodi usano gruppi di attributi \textbf{disgiunti}: sono due classi in una;
  \item per descrivere la classe serve una frase con molte virgole;
  \item una modifica a un requisito qualunque finisce sempre per toccare quella classe.
\end{itemize}}
\end{minipage}

% ============================================================
\section{Accoppiamento}

\dfn{Accoppiamento}{
  L'\term{accoppiamento} (\textit{coupling}) misura quanto una classe è \textbf{connessa} ad altre classi, quanto ne
  \textbf{conosce} i dettagli e quanto ne \textbf{dipende}.

  Due o più classi sono accoppiate quando è impossibile riusarne una senza riusare anche le altre. Se due classi dipendono
  strettamente e in molti dettagli l'una dall'altra, si dicono \textbf{fortemente accoppiate}; altrimenti
  \textbf{debolmente accoppiate}. Per avere codice di qualità bisogna puntare a un \textbf{basso accoppiamento}.
}

Il basso accoppiamento permette di:
\begin{itemize}
  \item capire il codice di una classe senza leggere i dettagli delle altre;
  \item modificare una classe senza che le modifiche si ripercuotano sulle altre;
  \item riusare facilmente una classe senza dover importare anche altre classi;
  \item testare facilmente una classe, indipendentemente dalle altre.
\end{itemize}
In una parola, migliora la \textbf{manutenibilità} del software.

Un certo grado di accoppiamento è però \textbf{inevitabile}: nel paradigma a oggetti il sistema funziona proprio perché gli
oggetti si scambiano messaggi, e per mandare un messaggio a un oggetto bisogna conoscerlo. L'accoppiamento va quindi
\textbf{limitato} dove possibile, non eliminato. La linea guida è assegnare le responsabilità delle classi in modo da
ottenere un accoppiamento basso; la Legge di Demetra, più avanti, ne è un'applicazione concreta.

\subsection{Le forme di accoppiamento}

Date due classi $X$ e $Y$, $X$ è accoppiata a $Y$ se si verifica almeno uno di questi casi. Accanto riportiamo come la
dipendenza si vede in un diagramma delle classi (Capitolo 7).

\Needspace{16\baselineskip}
\begin{center}
\renewcommand{\arraystretch}{1.6}
\begin{tabular}{@{}c p{6.4cm} c@{}}
  \toprule
  & \textbf{Caso} & \textbf{In UML} \\
  \midrule
  1 & $X$ ha un \textbf{attributo} di tipo $Y$ &
    \raisebox{-0.25\height}{\tikz{\node[classe1=0.8cm](x){X}; \node[classe1=0.8cm](y) at (2.4,0){Y}; \draw[navig](x)--(y);}} \\
  2 & $X$ ha un \textbf{metodo} che usa un elemento di tipo $Y$ (parametro, variabile locale, valore restituito\ldots) &
    \raisebox{-0.25\height}{\tikz{\node[classe1=0.8cm](x){X}; \node[classe1=0.8cm](y) at (2.4,0){Y}; \draw[dipend](x)--(y);}} \\
  3 & $X$ è una \textbf{sottoclasse} (anche indiretta) di $Y$ &
    \raisebox{-0.25\height}{\tikz{\node[classe1=0.8cm](x){X}; \node[classe1=0.8cm](y) at (2.4,0){Y}; \draw[gen](x)--(y);}} \\
  4 & $X$ \textbf{implementa} un'interfaccia di tipo $Y$ &
    \raisebox{-0.25\height}{\tikz{\node[classe1=0.8cm](x){X}; \node[classe1=0.8cm, font=\tiny\bfseries, align=center](y) at (2.4,0){«interface»\\Y}; \draw[realizza](x)--(y);}} \\
  \bottomrule
\end{tabular}
\end{center}

Il caso 3, l'\textbf{ereditarietà}, è quello che produce l'accoppiamento più forte. Una sottoclasse eredita attributi e
metodi della superclasse e spesso ne dipende anche nei dettagli di implementazione: una modifica alla superclasse può
rompere tutte le sottoclassi, anche quelle che non si sono mai occupate della parte modificata. L'interfaccia (caso 4),
invece, impegna solo a rispettare un \textbf{contratto}, senza portarsi dietro alcuna implementazione.

\subsection{Pensare per interfacce}

Le interfacce sono uno strumento centrale per ridurre l'accoppiamento, e la libreria standard di Java ne contiene centinaia
di esempi d'uso corretto (\code{List}, \code{Comparable}, \code{Runnable}\ldots). I loro vantaggi principali sono due.
\begin{itemize}
  \item \textbf{Fissano uno standard.} Un'interfaccia definisce un \textbf{contratto} che favorisce il riuso. Un oggetto che
        implementa un'interfaccia promette di rispettare quello standard; l'oggetto che lo usa si chiama
        \term{consumatore} (\textit{consumer}). L'interfaccia è un contratto fra un oggetto e i suoi consumatori.
  \item \textbf{Introducono un livello di astrazione} che rende i programmi più facili da capire: permettono di ragionare
        sul comportamento generale del codice senza entrare nei dettagli specifici.
\end{itemize}

Vediamo l'effetto di un'interfaccia su un esempio concreto.

\begin{esempio}{Il viaggiatore e l'automobile}
\Needspace{14\baselineskip}
\noindent\begin{minipage}[c]{0.5\textwidth}
\begin{lstlisting}[language=Java]
class Traveler {
  Car c = new Car();

  void startJourney() {
    c.move();
  }
}

class Car {
  void move() {
    // logica...
  }
}
\end{lstlisting}
\end{minipage}\hfill
\begin{minipage}[c]{0.46\textwidth}
\centering
\begin{tikzpicture}
  \umlclass[2.3cm]{t}{(0,0)}{Traveler}{c : Car}{startJourney()}
  \umlclass[1.8cm]{c}{(3.6,0)}{Car}{}{move()}
  \draw[navig] (t.east) -- (c.west);
\end{tikzpicture}
\end{minipage}

\code{Traveler} ha un attributo di tipo \code{Car}: è il caso 1, un accoppiamento forte.
\begin{itemize}
  \item Che cosa succede se vogliamo riusare \code{Traveler} con un mezzo di trasporto diverso da \code{Car}? Dobbiamo
        modificarlo.
  \item Che cosa succede se cambiano i metodi di \code{Car}? Dobbiamo modificare anche \code{Traveler}.
\end{itemize}

Introduciamo un'interfaccia \code{Vehicle}, che \code{Car} e \code{Bike} implementano.

\Needspace{18\baselineskip}
\noindent\begin{minipage}[c]{0.5\textwidth}
\begin{lstlisting}[language=Java]
interface Vehicle {
  void move();
}

class Car implements Vehicle {
  public void move() { /* logica */ }
}

class Bike implements Vehicle {
  public void move() { /* logica */ }
}

class Traveler {
  Vehicle v;

  public void setV(Vehicle v) {
    this.v = v;
  }
  void startJourney() {
    v.move();
  }
}
\end{lstlisting}
\end{minipage}\hfill
\begin{minipage}[c]{0.46\textwidth}
\centering
\begin{tikzpicture}
  \umlclass[2.3cm]{t}{(0,0)}{Traveler}{v : Vehicle}{setV(v : Vehicle)\\startJourney()}
  \node[classe1=2.0cm, font=\scriptsize, align=center] (v) at (3.8,0.35) {«interface»\\\textbf{Vehicle}};
  \node[font=\tiny, anchor=north] at (v.south) {move()};
  \umlclass[1.5cm]{c}{(2.8,-2.2)}{Car}{}{move()}
  \umlclass[1.5cm]{b}{(4.8,-2.2)}{Bike}{}{move()}
  \draw[navig] (t.east |- v.west) -- (v.west);
  \draw[realizza] (c.north) -- ($(v.south)+(-0.4,-0.3)$);
  \draw[realizza] (b.north) -- ($(v.south)+(0.4,-0.3)$);
\end{tikzpicture}

\vspace{4pt}
{\scriptsize \code{Traveler} non conosce più né \code{Car} né \code{Bike}.}
\end{minipage}

L'interfaccia \code{Vehicle} ha \textbf{spezzato} l'accoppiamento fra \code{Traveler} e \code{Car}. Le conseguenze:
\begin{itemize}
  \item \code{Traveler} ora è accoppiato a un'\textbf{interfaccia}, non a una classe concreta;
  \item tutte le implementazioni di \code{Vehicle} funzionano con \code{Traveler} senza modificarne il codice;
  \item si potranno usare in futuro implementazioni di \code{Vehicle} che oggi non esistono ancora (\textbf{anticipare il
        cambiamento});
  \item serve però un'\textbf{entità esterna} che fornisca (\textit{inietti}) l'implementazione concreta: qui il metodo
        \code{setV}, chiamato ad esempio così:
\begin{lstlisting}[language=Java]
Traveler t = new Traveler();
t.setV(new Bike());      // chi crea Traveler sceglie il mezzo
t.startJourney();
\end{lstlisting}
  \item la complessità del codice è aumentata: c'è un tipo in più da conoscere.
\end{itemize}
\end{esempio}

\warning{{Iniezione con il costruttore}}{
  Nella versione delle slide, se qualcuno chiama \code{startJourney()} senza aver prima chiamato \code{setV(...)},
  l'attributo \code{v} vale \code{null} e il programma termina con una \code{NullPointerException}. Quando la dipendenza
  è \textbf{obbligatoria}, conviene passarla nel \textbf{costruttore}: \code{public Traveler(Vehicle v) \{ this.v = v; \}}.
  Così non può esistere un \code{Traveler} senza mezzo di trasporto. Questa tecnica, con cui chi usa una classe le fornisce
  dall'esterno gli oggetti di cui ha bisogno, si chiama \term{dependency injection}. Applicata a \code{ReportPrinter}
  dell'esempio sul report, risolverebbe il problema di accoppiamento segnalato lì.
}

% ============================================================
\section{Un esempio completo: il ragazzo dei giornali}

L'esempio classico del \textit{Paperboy} (il ragazzo che consegna i giornali e passa a riscuotere) mostra un caso di
accoppiamento eccessivo molto comune, e come correggerlo.

\Needspace{24\baselineskip}
\noindent\begin{minipage}[t]{0.49\textwidth}
\begin{lstlisting}[language=Java]
public class Customer {
  private String firstName;
  private String lastName;
  private Wallet myWallet;

  public String getFirstName() {
    return firstName;
  }
  public String getLastName() {
    return lastName;
  }
  public Wallet getWallet() {
    return myWallet;
  }
}
\end{lstlisting}
\end{minipage}\hfill
\begin{minipage}[t]{0.49\textwidth}
\begin{lstlisting}[language=Java]
public class Wallet {
  private float value;

  public float getTotalMoney() {
    return value;
  }
  public void setTotalMoney(float newValue) {
    value = newValue;
  }
  public void addMoney(float deposit) {
    value += deposit;
  }
  public void subtractMoney(float debit) {
    value -= debit;
  }
}
\end{lstlisting}
\end{minipage}

Ecco il codice di un metodo della classe \code{Paperboy}, che riscuote due dollari.

\Needspace{10\baselineskip}
\begin{lstlisting}[language=Java]
float payment = 2.00f;                     // "voglio i miei due dollari!"
Wallet theWallet = myCustomer.getWallet();
if (theWallet.getTotalMoney() > payment) {
  theWallet.subtractMoney(payment);
} else {
  // ripassare piu' tardi a riscuotere
}
\end{lstlisting}

\subsection{Perché è sbagliato}

\begin{enumerate}
  \item \textbf{Non rispecchia il mondo reale.} Tradotto in linguaggio naturale, il codice dice che quando il ragazzo
        passa a riscuotere, il cliente gli lascia \textbf{prendere il portafoglio} dalla tasca e sfilarne due dollari,
        fidandosi che prenda solo quanto dovuto. Se in futuro \code{Wallet} contenesse anche le carte di credito, il
        ragazzo avrebbe accesso pure a quelle. Il problema di fondo è che il \code{Paperboy} è \textbf{esposto a più
        informazioni di quante gliene servano}.
  \item \textbf{Accoppia tre classi.} \code{Paperboy} ora sa che il cliente ha un portafoglio e può manipolarlo: per
        compilarlo servono sia \code{Customer} sia \code{Wallet}. Le tre classi sono \textbf{fortemente accoppiate}, e una
        modifica a \code{Wallet} può costringere a modificare le altre due.
  \Needspace{16\baselineskip}
  \item \textbf{È fragile.} Supponiamo che a un cliente rubino il portafoglio, e che un collega decida di modellare il
        furto impostando il portafoglio a \code{null}:
\begin{lstlisting}[language=Java]
victim.setWallet(null);
\end{lstlisting}
        È una scelta ragionevole: né la documentazione né il codice impongono che il
        portafoglio esista. Ma il codice del \code{Paperboy} \textbf{dà per scontato} che ci sia, e chiamando un metodo su
        \code{null} ottiene una \code{NullPointerException}. Si potrebbe controllare \code{null} prima di ogni chiamata, ma il
        codice del ragazzo diventerebbe ancora più confuso: \textit{«se il cliente ha un portafoglio, guarda quanto c'è
        dentro; se può pagarmi, prendi i soldi»}.
\end{enumerate}

\begin{center}
\begin{tikzpicture}
  \node[classe1=2.0cm] (p) at (0,0) {Paperboy};
  \node[classe1=2.0cm] (c) at (3.4,0) {Customer};
  \node[classe1=2.0cm] (w) at (6.8,0) {Wallet};
  \draw[dipend] (p) -- (c);
  \draw[dipend] (c) -- (w);
  \draw[dipend, draw=swred, line width=0.9pt] (p.north) to[bend left=25]
    node[note, above, text=swred]{il ragazzo conosce e usa il portafoglio} (w.north);
  \node[note, font=\scriptsize\bfseries] at (3.4,-0.75) {prima};
  \begin{scope}[shift={(9.4,0)}]
    \node[classe1=2.0cm] (p2) at (0,0) {Paperboy};
    \node[classe1=2.0cm] (c2) at (3.4,0) {Customer};
    \node[classe1=2.0cm] (w2) at (3.4,1.2) {Wallet};
    \draw[dipend] (p2) -- (c2);
    \draw[dipend] (c2) -- (w2);
    \node[note, font=\scriptsize\bfseries] at (1.7,-0.75) {dopo};
  \end{scope}
\end{tikzpicture}
\end{center}

\subsection{La versione migliorata}

\code{Customer} non offre più \code{getWallet()}, ma un metodo \code{getPayment()}: il ragazzo \textbf{chiede} un pagamento,
e il cliente decide come farlo.

\Needspace{16\baselineskip}
\noindent\begin{minipage}[t]{0.52\textwidth}
\begin{lstlisting}[language=Java]
public class Customer {
  private String firstName;
  private String lastName;
  private Wallet myWallet;

  // getFirstName(), getLastName() come prima

  public float getPayment(float bill) {
    if (myWallet != null
        && myWallet.getTotalMoney() >= bill) {
      myWallet.subtractMoney(bill);
      return bill;
    }
    return 0;
  }
}
\end{lstlisting}
\end{minipage}\hfill
\begin{minipage}[t]{0.45\textwidth}
\begin{lstlisting}[language=Java]
// in un metodo di Paperboy
float payment = 2.00f;
float paidAmount =
    myCustomer.getPayment(payment);
if (paidAmount == payment) {
  // ringrazia e lascia la ricevuta
} else {
  // ripassa piu' tardi
}
\end{lstlisting}
\end{minipage}

\warning{Il codice delle slide va corretto}{
  Il metodo \code{getPayment} delle slide non compila:
  \begin{itemize}
    \item usa le variabili \code{theWallet} e \code{payment}, che appartengono al codice del \code{Paperboy}: dentro
          \code{Customer} vanno usati l'attributo \code{myWallet} e il parametro \code{bill};
    \item se il portafoglio manca o non contiene abbastanza soldi, il metodo termina \textbf{senza restituire nulla}, e
          Java segnala \textit{missing return statement}: serve un \code{return 0} finale.
  \end{itemize}
  Inoltre, sia nelle slide sia nel codice originale, la condizione \code{getTotalMoney() > payment} impedisce di pagare a
  un cliente che ha \textbf{esattamente} due dollari: la condizione corretta è \code{>=}. Infine, in un'applicazione reale
  le somme di denaro non si rappresentano con \code{float}, che introduce errori di arrotondamento, ma con
  \code{BigDecimal} o con un intero di centesimi.
}

\subsection{Perché è meglio}

\begin{enumerate}
  \item \textbf{Modella meglio la realtà.} Il codice del \code{Paperboy} ora \textit{chiede} al cliente un pagamento, e non
        ha accesso diretto al portafoglio.
  \item \textbf{Isola i cambiamenti.} La classe \code{Wallet} può cambiare e il ragazzo ne è completamente isolato. Se
        cambiasse l'interfaccia di \code{Wallet}, andrebbe aggiornato solo \code{Customer}; finché l'interfaccia di
        \code{Customer} resta la stessa, nessun suo cliente si accorgerà che ha un portafoglio nuovo. Il codice è più
        manutenibile perché le modifiche non si \textbf{propagano a catena} in un grande progetto.
  \item \textbf{Libera l'implementazione.} È forse la ragione più «a oggetti»: ora siamo liberi di cambiare
        l'implementazione di \code{getPayment()}. Nella realtà, quando passa il ragazzo, il cliente potrebbe prendere i due
        dollari da un barattolo di monete, cercarli fra i cuscini del divano o farseli prestare dal coinquilino. Tutto questo
        è \textbf{logica di business} che non riguarda il ragazzo: può stare dentro \code{getPayment()} e cambiare in futuro
        senza toccare il codice del \code{Paperboy}.
\end{enumerate}

\info{{Ma Customer non è diventato più complesso?}}{
  Un'obiezione legittima: abbiamo reso \code{Customer} più complesso. Ma nella prima versione era il \code{Paperboy} a dover
  conoscere e manipolare il \code{Wallet}. Se \code{Customer} è più complesso, perché i suoi clienti ora sono \textbf{meno}
  complessi? Perché la complessità è sempre la stessa: è solo \textbf{contenuta} meglio nell'oggetto giusto ed esposta in
  modo sano. Prima esisteva nel codice dei clienti, magari ripetuta in più punti; ora compare \textbf{una volta sola}, e si
  mantiene nello stesso posto in cui avvengono i cambiamenti che potrebbero romperla.
}

% ============================================================
\section{La Legge di Demetra}

La \term{Legge di Demetra} (\textit{Law of Demeter}) è un'altra linea guida fondamentale per ottenere un basso
accoppiamento. Prende il nome dal progetto Demeter, uno dei primi grandi sistemi a oggetti, sviluppato a Boston alla fine degli
anni Ottanta (le slide indicano la Boston University, ma il gruppo di Karl Lieberherr che formulò la legge lavorava alla
Northeastern University).

L'idea di fondo è spingere l'\textbf{information hiding} al massimo: niente catene del tipo
\code{foo.getA().getB().getC()}. Più lungo è il percorso di chiamate che il programma attraversa, più il programma è fragile
ai cambiamenti: ogni classe lungo la catena è una classe da cui dipendiamo.

\dfn{Legge di Demetra}{
  Un metodo $M$ di un oggetto $O$ dovrebbe inviare messaggi \textbf{solo} a:
  \begin{enumerate}
    \item l'oggetto stesso (\code{this});
    \item i \textbf{parametri} di $M$;
    \item gli oggetti \textbf{creati} all'interno di $M$;
    \item i \textbf{componenti diretti} di $O$, cioè i suoi attributi.
  \end{enumerate}
  In breve: \textit{parla solo con i tuoi amici più stretti, non con gli amici degli amici}.
}

\begin{esempio}{Chiamate permesse e vietate}
\Needspace{16\baselineskip}
\begin{lstlisting}[language=Java]
class Driver {
  private Car car;                          // componente diretto

  void drive(Road road) {                   // parametro
    this.checkMirrors();                    // OK (1): this
    road.getSpeedLimit();                   // OK (2): parametro
    Route r = new Route(); r.compute();     // OK (3): creato qui
    car.increaseSpeed();                    // OK (4): componente diretto
    car.getEngine().getCarburator()
       .openFuelValve();                    // NO: oggetti estranei
  }
}
\end{lstlisting}
L'ultima chiamata attraversa \code{Car}, \code{Engine} e \code{Carburator}: \code{Driver} dipende così da tre classi invece
che da una, e da come sono collegate fra loro.
\end{esempio}

\Needspace{14\baselineskip}
\subsection{Com'è fatta una violazione}

In figura un \textit{client} raggiunge un fornitore indiretto passando attraverso un intermediario. Le due scritture
nella nota sono equivalenti: usare un \textit{getter} invece dell'accesso diretto all'attributo \textbf{non} rimedia alla
violazione, perché il client continua a conoscere la struttura interna dell'intermediario.

\begin{center}
\begin{tikzpicture}
  \umlclass[2.6cm]{ip}{(0,0)}{Indirect Provider}{}{doSomething()}
  \umlclass[2.6cm]{im}{(5.0,0)}{Intermediate Provider}{+provider}{getProvider()}
  \node[classe1=2.2cm] (cl) at (10.0,0.25) {Indirect Client};
  \draw[assoc] (ip.east) -- node[above, font=\scriptsize\itshape]{provider} (im.west);
  \draw[assoc] (im.east |- cl.west) -- node[above, font=\scriptsize\itshape]{intermediate} (cl.west);
  \node[notauml, font=\scriptsize\ttfamily, anchor=north] (n) at (5.0,-1.5)
    {intermediate.provider.doSomething()\\ \textrm{\textit{oppure}}\\ intermediate.getProvider().doSomething()};
  \draw[dashed, draw=black!60] (n.north east) -- (cl.south);
  \draw[swred, line width=2.2pt, opacity=0.8] (n.center) circle (0.62);
  \draw[swred, line width=2.2pt, opacity=0.8] ($(n.center)+(-0.44,0.44)$) -- ($(n.center)+(0.44,-0.44)$);
\end{tikzpicture}
\end{center}

\subsection{Eliminare il codice di navigazione}

La correzione consiste nel \textbf{spostare il comportamento} verso l'oggetto che possiede i dati, un passo alla volta.
L'esempio (da Oscar Nierstrasz) è un'auto che accelera aprendo la valvola del carburatore.

\begin{figure}[H]
\centering
\begin{tikzpicture}[
    cod/.style={notauml, font=\scriptsize\ttfamily},
    passo/.style={chip=swgray, font=\scriptsize\bfseries},
    rif/.style={dashed, draw=black!55, line width=0.5pt}]
  % --- prima ---
  \node[passo, anchor=east] at (-1.55,0) {Prima};
  \umlclass[2.6cm]{c0}{(0,0)}{Carburator}{+fuelValveOpen}{}
  \umlclass[2.6cm]{e0}{(4.2,0)}{Engine}{+carburator}{}
  \umlclass[2.6cm]{a0}{(8.4,0)}{Car}{-engine}{+increaseSpeed()}
  \draw[assoc] (c0.east) -- (e0.west);
  \draw[assoc] (e0.east) -- (a0.west);
  \node[cod, anchor=north east, text=swred] (n0) at ($(a0.south east)+(0,-0.3)$)
    {engine.carburator.fuelValveOpen = true};
  \draw[rif] (n0.north -| a0.south) -- (a0.south);
  % --- passo 1 ---
  \node[passo, anchor=east] at (-1.55,-2.9) {Passo 1};
  \umlclass[2.6cm]{c1}{(0,-2.9)}{Carburator}{+fuelValveOpen}{}
  \umlclass[2.6cm]{e1}{(4.2,-2.9)}{Engine}{-carburator}{\textcolor{swred}{+speedUp()}}
  \umlclass[2.6cm]{a1}{(8.4,-2.9)}{Car}{-engine}{+increaseSpeed()}
  \draw[assoc] (c1.east) -- (e1.west);
  \draw[assoc] (e1.east) -- (a1.west);
  \node[cod, anchor=north] (n1a) at ($(e1.south)+(0,-0.3)$) {carburator.fuelValveOpen = true};
  \node[cod, anchor=north] (n1b) at ($(a1.south)+(0.3,-0.3)$) {engine.speedUp()};
  \draw[rif] (n1a.north) -- (e1.south);
  \draw[rif] (n1b.north -| a1.south) -- (a1.south);
  % --- passo 2 ---
  \node[passo, anchor=east] at (-1.55,-5.8) {Passo 2};
  \umlclass[2.6cm]{c2}{(0,-5.8)}{Carburator}{-fuelValveOpen}{\textcolor{swred}{+openFuelValve()}}
  \umlclass[2.6cm]{e2}{(4.2,-5.8)}{Engine}{-carburator}{+speedUp()}
  \umlclass[2.6cm]{a2}{(8.4,-5.8)}{Car}{-engine}{+increaseSpeed()}
  \draw[assoc] (c2.east) -- (e2.west);
  \draw[assoc] (e2.east) -- (a2.west);
  \node[cod, anchor=north] (n2a) at ($(c2.south)+(0,-0.3)$) {fuelValveOpen = true};
  \node[cod, anchor=north] (n2b) at ($(e2.south)+(0,-0.3)$) {carburator.openFuelValve()};
  \node[cod, anchor=north] (n2c) at ($(a2.south)+(0.3,-0.3)$) {engine.speedUp()};
  \draw[rif] (n2a.north) -- (c2.south);
  \draw[rif] (n2b.north) -- (e2.south);
  \draw[rif] (n2c.north -| a2.south) -- (a2.south);
\end{tikzpicture}
\caption{Eliminare il codice di navigazione. A ogni passo un oggetto smette di frugare nella struttura di un altro e gli
         chiede invece un servizio (in rosso i metodi aggiunti; nelle note il codice dei metodi).}
\end{figure}

\begin{itemize}
  \item \textbf{Prima}: \code{Car} attraversa \code{Engine} e scrive direttamente l'attributo del \code{Carburator}. Conosce
        la struttura interna di due classi.
  \item \textbf{Passo 1}: \code{Engine} offre il servizio \code{speedUp()}. \code{Car} parla solo con il suo componente
        diretto, l'\code{engine}; l'attributo \code{carburator} può diventare privato.
  \item \textbf{Passo 2}: anche \code{Carburator} offre un servizio, \code{openFuelValve()}, e il suo attributo diventa
        privato. Ora ogni classe conosce solo i propri vicini immediati.
\end{itemize}

\info{{Il Paperboy e Demetra}}{
  La prima versione del \code{Paperboy} violava proprio la Legge di Demetra: \code{myCustomer.getWallet().subtractMoney(...)}
  manda un messaggio al \code{Wallet}, che non è né un parametro né un componente del ragazzo. La correzione è la stessa
  dell'auto: spostare il comportamento (\code{getPayment()}) nell'oggetto che possiede i dati.

  Attenzione a non ridurre la regola al conteggio dei punti: \code{builder.append("a").append("b")} o una catena di
  operazioni su uno \textit{stream} non violano la Legge di Demetra, perché ogni chiamata restituisce lo stesso oggetto o un
  oggetto dello stesso tipo, e non si naviga nella struttura interna di oggetti estranei.
}

% ============================================================
\section{Progettazione guidata dalle responsabilità}

La \term{progettazione guidata dalle responsabilità} (\textit{responsibility-driven design}) ragiona sul sistema in termini
di tre concetti: \textbf{classe}, \textbf{responsabilità} e \textbf{collaborazione}. Mette insieme le regole viste finora.

\begin{center}
\begin{tikzpicture}
  \node[card=swgreen, text width=4.6cm, minimum height=1.5cm] (a) at (0,0)
    {ogni classe ha responsabilità e ruoli \textbf{chiaramente definiti}};
  \node[card=swblue, text width=4.6cm, minimum height=1.5cm] (b) at (5.3,0)
    {la classe che \textbf{possiede i dati} è responsabile di elaborarli};
  \node[card=q7, text width=4.6cm, minimum height=1.5cm] (c) at (10.6,0)
    {ogni classe è responsabile della \textbf{gestione dei propri dati}};
  \node[chip=swgreen, anchor=north] at ($(a.south)+(0,-0.15)$) {alta coesione};
  \node[chip=swblue, anchor=north] at ($(b.south)+(0,-0.15)$) {basso accoppiamento};
  \node[chip=q7, anchor=north] at ($(c.south)+(0,-0.15)$) {incapsulamento};
\end{tikzpicture}
\end{center}

Lo strumento che facilita questo modo di ragionare sono le \textbf{schede CRC}, già incontrate nel Capitolo 7 per
identificare le classi del modello di dominio.

\subsection{Le schede CRC}

Le schede CRC (\textit{Class--Responsibility--Collaboration}) sono state introdotte da Kent Beck e Ward Cunningham per
insegnare la progettazione a oggetti. A ogni classe del diagramma delle classi di progetto si associa una scheda che
riporta:
\begin{itemize}
  \item il \textbf{nome} della classe;
  \item le sue \textbf{responsabilità}: le conoscenze che la classe mantiene o i servizi che offre;
  \item i suoi \textbf{collaboratori}: le altre classi di cui la classe deve sfruttare conoscenze o servizi per assolvere
        le proprie responsabilità.
\end{itemize}

\Needspace{20\baselineskip}
\noindent\begin{minipage}[c]{0.4\textwidth}
\centering
\begin{tikzpicture}[font=\scriptsize]
  \draw[line width=0.8pt, fill=swlight] (0,0) rectangle (5.4,4.6);
  \draw (0,4.1) -- (5.4,4.1);
  \draw[line width=1.2pt] (0,3.6) -- (5.4,3.6);
  \draw (0,3.1) -- (5.4,3.1);
  \draw[line width=1.2pt] (0,2.6) -- (5.4,2.6);
  \draw (3.2,0) -- (3.2,2.6);
  \foreach \y in {0.4,0.8,...,2.2} \draw[swgray!50] (0,\y) -- (5.4,\y);
  \node[anchor=west] at (0.08,4.35) {\textbf{Nome della classe}};
  \node[anchor=west] at (0.08,3.85) {Superclassi};
  \node[anchor=west] at (0.08,3.35) {Sottoclassi};
  \node[anchor=west] at (0.08,2.82) {Responsabilità};
  \node[anchor=west] at (3.28,2.82) {Collaboratori};
\end{tikzpicture}

{\scriptsize La scheda CRC completa delle slide.}
\end{minipage}\hfill
\begin{minipage}[c]{0.56\textwidth}
Rispetto alla versione del Capitolo 7 questa scheda riporta anche \textbf{superclassi} e \textbf{sottoclassi}: in
progettazione la gerarchia è decisa, e va tenuta sotto controllo perché l'ereditarietà è la forma di accoppiamento più forte.

Le schede CRC hanno due conseguenze pratiche, che ne fanno uno strumento di \textbf{diagnosi}:
\begin{itemize}
  \item se non si riescono a specificare con chiarezza le responsabilità di una classe, c'è un \textbf{errore di progetto}
        (bassa coesione);
  \item se una classe ha bisogno di \textbf{troppi collaboratori}, c'è un potenziale errore di progetto (alto
        accoppiamento).
\end{itemize}
\end{minipage}

\begin{esempio}{Le schede CRC del Paperboy}
Ecco le schede della versione migliorata. Si noti che \code{Wallet} non compare fra i collaboratori del \code{Paperboy}:
la scheda rende visibile a colpo d'occhio la riduzione dell'accoppiamento.

\begin{center}
\begin{tikzpicture}
  \crccard{(0,0)}{Paperboy}{consegna i giornali\newline riscuote il pagamento}{Customer}
  \crccard{(6.1,0)}{Customer}{conosce nome e cognome\newline paga un conto (\code{getPayment})}{Wallet}
  \crccard{(3.05,-1.9)}{Wallet}{conosce la somma contenuta\newline aggiunge e toglie denaro}{---}
\end{tikzpicture}
\end{center}
\end{esempio}

% ============================================================
\section{Esercizi}

\begin{esercizio}{Troppe responsabilità}
La classe seguente fa parte di un sistema di segreteria universitaria. Individua le sue responsabilità, spiega quali motivi
di cambiamento la rendono fragile e proponi una scomposizione che rispetti il Single Responsibility Principle.

\begin{lstlisting}[language=Java]
public class Studente {
  private String matricola, nome, email;
  private List<Esame> esami;

  public double calcolaMedia() { ... }
  public void salvaSuDatabase(Connection db) { ... }
  public void stampaLibretto() { ... }
  public void inviaEmailConvocazione(String testo) { ... }
}
\end{lstlisting}

\textbf{Soluzione.} Le responsabilità sono quattro, e ognuna è un motivo di cambiamento diverso:
\begin{itemize}
  \item i \textbf{dati} dello studente e il calcolo della media (cambia se cambiano le regole di calcolo, ad esempio la
        media pesata sui crediti);
  \item la \textbf{persistenza} (cambia se cambia il DBMS o lo schema delle tabelle);
  \item la \textbf{presentazione} del libretto (cambia se si vuole un PDF invece della stampa a video);
  \item la \textbf{comunicazione} (cambia se si passa dalle e-mail alle notifiche dell'app).
\end{itemize}
Una scomposizione possibile: \code{Studente} tiene i dati e \code{calcolaMedia()}, perché la media si calcola dagli esami,
che sono dati suoi; \code{StudenteRepository} si occupa di salvare e caricare; \code{LibrettoPrinter} della stampa;
\code{ServizioNotifiche} dell'invio dei messaggi.

\begin{center}
\begin{tikzpicture}
  \umlclass[2.9cm]{s}{(0,0)}{Studente}{matricola, nome, email\\esami}{calcolaMedia()}
  \umlclass[3.4cm]{r}{(5.6,1.5)}{StudenteRepository}{}{salva(s : Studente)\\carica(matr) : Studente}
  \umlclass[3.4cm]{p}{(5.6,0)}{LibrettoPrinter}{}{stampa(s : Studente)}
  \umlclass[3.4cm]{n}{(5.6,-1.4)}{ServizioNotifiche}{}{invia(s : Studente, testo)}
  \draw[dipend] (r.west) -- (s.east);
  \draw[dipend] (p.west) -- (s.east);
  \draw[dipend] (n.west) -- (s.east);
\end{tikzpicture}
\end{center}
Le tre classi nuove dipendono da \code{Studente}, mentre \code{Studente} non dipende da nessuna di loro: si può cambiare il
database o il canale di notifica senza toccarla.
\end{esercizio}

\begin{esercizio}{Riconoscere l'accoppiamento}
Elenca tutte le classi a cui è accoppiata \code{Ordine}, indicando per ciascuna il tipo di accoppiamento. Quale conviene
ridurre per primo?

\begin{lstlisting}[language=Java]
public class Ordine extends Documento implements Stampabile {
  private Cliente cliente;

  public void aggiungi(Prodotto p) {
    Sconto s = new Sconto();
    // ...
  }
}
\end{lstlisting}

\textbf{Soluzione.}

\begin{center}
\renewcommand{\arraystretch}{1.3}
\begin{tabular}{@{}lll@{}}
  \toprule
  \textbf{Classe} & \textbf{Caso} & \textbf{Perché} \\
  \midrule
  \code{Documento}  & 3, sottoclasse    & \code{Ordine extends Documento} \\
  \code{Stampabile} & 4, interfaccia    & \code{Ordine implements Stampabile} \\
  \code{Cliente}    & 1, attributo      & \code{private Cliente cliente} \\
  \code{Prodotto}   & 2, metodo         & parametro di \code{aggiungi} \\
  \code{Sconto}     & 2, metodo         & variabile locale creata in \code{aggiungi} \\
  \bottomrule
\end{tabular}
\end{center}
L'accoppiamento più forte è quello con \code{Documento}: conviene chiedersi se un ordine \textit{è davvero} un documento o
se piuttosto \textit{ha} un documento associato, e in quel caso sostituire l'ereditarietà con un attributo (composizione).
Anche \code{new Sconto()} lega \code{Ordine} a una classe concreta: se le politiche di sconto possono cambiare, è meglio
ricevere uno \code{Sconto} (o un'interfaccia \code{PoliticaSconto}) dall'esterno, come per \code{Vehicle}.
\end{esercizio}

\begin{esercizio}{Applicare la Legge di Demetra}
In un sistema di e-commerce, la classe \code{Spedizione} contiene il metodo seguente. Spiega perché viola la Legge di
Demetra e correggilo.

\begin{lstlisting}[language=Java]
public double costo(Ordine ordine) {
  String citta = ordine.getCliente().getIndirizzo().getCitta();
  return tariffario.tariffaPer(citta);
}
\end{lstlisting}

\textbf{Soluzione.} \code{ordine} è un parametro, quindi chiamare \code{ordine.getCliente()} è lecito; ma il risultato è un
\code{Cliente}, che non è né un parametro né un componente di \code{Spedizione}, e ancora meno lo è l'\code{Indirizzo}.
\code{Spedizione} dipende così da \code{Ordine}, \code{Cliente} e \code{Indirizzo}, e dal modo in cui sono collegati: se un
giorno l'indirizzo di consegna diventasse diverso da quello del cliente, il metodo darebbe un risultato sbagliato.

La correzione sposta la responsabilità verso chi ha i dati: è l'\code{Ordine} a sapere dove va consegnato.
\begin{lstlisting}[language=Java]
// in Ordine
public String getCittaConsegna() {
  return cliente.getCittaConsegna();   // Cliente delega al suo Indirizzo
}

// in Spedizione
public double costo(Ordine ordine) {
  return tariffario.tariffaPer(ordine.getCittaConsegna());
}
\end{lstlisting}
Un'alternativa altrettanto valida è far ricevere a \code{costo} direttamente la città, o l'\code{Indirizzo}, come
parametro.
\end{esercizio}

\begin{esercizio}{Obiettivi di progetto e compromessi}
Una regione commissiona un'app per prenotare le vaccinazioni. Deve essere online fra \textbf{due mesi}, all'inizio della
campagna; nei primi giorni si attendono centinaia di migliaia di accessi, in gran parte da persone anziane. L'app tratta
dati sanitari. Indica i tre obiettivi di progetto prioritari e almeno un compromesso da dichiarare.

\textbf{Soluzione.} Una risposta ragionevole:
\begin{enumerate}
  \item \textbf{disponibilità e prestazioni sotto carico}: il picco iniziale è certo, e un sistema che non risponde il primo
        giorno fallisce la sua missione;
  \item \textbf{sicurezza}: i dati sanitari sono dati particolari secondo il GDPR;
  \item \textbf{usabilità}: gli utenti principali hanno poca dimestichezza con la tecnologia.
\end{enumerate}
Il compromesso principale è \textbf{tempi di rilascio \textit{vs} funzionalità}: la data non si può spostare, quindi si
rilascia una prima versione con le sole funzioni essenziali (prenotare, spostare, annullare) e si rimandano le altre
(ad esempio lo storico delle vaccinazioni). Va dichiarato anche il compromesso \textbf{costo \textit{vs} affidabilità}:
servono server ridondanti dimensionati per il picco. È invece sconsigliato puntare sul compromesso tempi/personale a
progetto avviato, per la legge di Brooks.
\end{esercizio}

\begin{esercizio}{Schede CRC per una biblioteca}
Un utente prende in prestito una copia di un libro per 30 giorni. Il sistema calcola la data di restituzione e, in caso di
ritardo, una multa di 0,50 euro al giorno. Un utente non può avere più di 5 prestiti attivi. Scrivi le schede CRC di
\code{Utente}, \code{Copia} e \code{Prestito}.

\textbf{Soluzione.}
\begin{center}
\begin{tikzpicture}
  \crccard{(0,0)}{Prestito}{conosce copia, utente, data di inizio\newline calcola la data di scadenza\newline calcola la
    multa}{Copia\newline Utente}
  \crccard{(6.1,0)}{Copia}{conosce il libro e la collocazione\newline sa se è disponibile\newline si segna in prestito o
    restituita}{Libro}
  \crccard{(3.05,-2.3)}{Utente}{conosce i propri dati\newline conosce i prestiti attivi\newline sa se può fare un nuovo
    prestito (max 5)}{Prestito}
\end{tikzpicture}
\end{center}
La multa la calcola \code{Prestito}, perché è lui a possedere le date: è la regola «la classe che possiede i dati li
elabora». Allo stesso modo il controllo sul limite di 5 prestiti sta in \code{Utente}, che conosce i propri prestiti
attivi. Nessuna scheda ha più di due collaboratori: è un buon segnale di accoppiamento contenuto.
\end{esercizio}

% ============================================================
\gbox{In sintesi}{azure-gradient-6!80!black}{
\begin{itemize}
  \item Il \textbf{System Design} trasforma il modello di analisi nel modello di progetto, scritto per i
        \textbf{programmatori}: obiettivi di progetto, architettura in sottosistemi, strategie di costruzione. L'\textbf{Object
        Design} ne specifica poi i dettagli.
  \item Gli \textbf{obiettivi di progetto} derivano soprattutto dai requisiti non funzionali, vanno \textbf{formalizzati} e
        ordinati per priorità, perché si possono ottimizzare solo pochi criteri: memoria/velocità, tempi/funzionalità,
        tempi/qualità, tempi/personale sono i compromessi tipici.
  \item Si \textbf{scompone} perché $C(P_1+P_2) > C(P_1)+C(P_2)$; ogni sottosistema offre i suoi servizi tramite
        un'\textbf{interfaccia}. Il numero di componenti ottimale bilancia costo di sviluppo e costo di integrazione.
  \item L'obiettivo è l'\textbf{indipendenza funzionale}: \textbf{alta coesione} (una sola responsabilità, un solo motivo
        per cambiare) e \textbf{basso accoppiamento}.
  \item L'accoppiamento nasce da attributi, parametri e variabili, \textbf{ereditarietà} (il più forte) e interfacce; le
        \textbf{interfacce} con l'iniezione delle dipendenze lo riducono.
  \item La \textbf{Legge di Demetra}: un metodo parla solo con \code{this}, i parametri, gli oggetti che crea e i componenti
        diretti. Le catene \code{a.getB().getC()} si eliminano spostando il comportamento verso chi possiede i dati.
  \item Nella \textbf{progettazione guidata dalle responsabilità} le \textbf{schede CRC} servono da diagnosi: responsabilità
        poco chiare indicano bassa coesione, troppi collaboratori un accoppiamento alto.
\end{itemize}
}
